\section{模型评价与结论}
\label{sec:conclusion}

\subsection{主要结论}
\begin{enumerate}
  \item 建立了可复算的全量质量评价链条：22 维质量信号经方向统一、稳健归一化、CRITIC 客观赋权与加权距离 TOPSIS，输出 $261{,}086$ 条唯一记录的样本级质量分，并聚合至 7 个质量域与语料级。$A_1$ 语料级均分 $60.95$，全量语料 $54.23$，下降 $6.71$ 分的约 $99.8\%$ 由领域构成变化（github 占比 $19.52\%\to78.04\%$）解释；在两个有扩展数据的域上，域内质量变化不超过 $0.36$ 分。
  \item 提出了完全可补偿评分与有限补偿消解的递进式冲突处理框架：以五维相对分位阈值识别候选冲突（$A_1$ 冲突率 $29.52\%$，留出集 $29.47\%$），以闭式有限补偿模型（式~\ref{eq:q12-closed}）给出有界消解规则；扩展集与抽样集在 arxiv、github 两域上结论一致。
  \item 配比建模采用单纯形感知（ilr 坐标）的同尺度分层方案，并证明只含配比信息的跨尺度绝对损失预测在重复设计下结构不可辨识，据此将跨尺度结论限定为相对排序，避免不可辨识的外推。
  \item 设计了以损失为统一桥梁的四问接口：质量分 $Q_g$ 经 A16 映射进入广义标度律，弹性进入算力配置的边际比较，C6 桥接把损失翻译为 Benchmark 得分；半合成数据全程分层标注。
\end{enumerate}

\subsection{优点、局限与改进方向}
模型的优点在于全链路可追溯（每个数字对应 \texttt{run\_id}）、冲突消解有闭式解且参数有明确含义（$\lambda$ 即补偿上界）、对不可辨识的外推给出了排除性证明而非勉强点预测。局限在于质量评分尚缺 A18 外部锚定，冲突消解的两个建模自由度（阈值、五维分组）未做全组合稳健性扫描，问题二至四的结果尚未纳入。改进方向：以 A18 与 83 条标注清单完成外部效度闭环；对问题二在 $Q_g$ 接口上完成广义标度律的正式估计；问题三在标度律参数不确定性下给出预算分档的结构性转移判据。
